%% gene_family_detection.tex -- shared sub-recipe. Produces the family table and
%% the ortholog designation that the analysis recipes consume.

\begin{recipe}[label=sub:gene-family-detection]{Gene Family Detection}

\recipepurpose{Group genes across species into families, and decide which
  members of each family count as orthologs. Tensor Omics consumes both and
  produces neither: the grouping fixes what every later comparison is made
  \emph{within}, and the ortholog designation fixes what divergence is measured
  \emph{against}.}

\nomaintainer

\begin{requiredinputs}
  \item One protein FASTA per species, reduced to primary transcripts.
  \item Gene identifiers that match those in your expression table exactly ---
        the two are joined by string equality, nothing else.
  \item OrthoFinder, installed and run separately. It is not part of Tensor
        Omics and this recipe does not wrap it.
\end{requiredinputs}

\begin{prerequisites}
  \item None. Everything else in this manual depends on this step.
\end{prerequisites}

\begin{workflow}
  \step{Reduce each proteome to one protein per gene, the primary transcript.
        Isoforms are the single most common cause of inflated families.}
  \step{Run OrthoFinder on the directory of proteomes, following its own
        documentation. Nothing about this step is Tensor Omics specific.}
  \step{Take \file{Orthogroups/Orthogroups.tsv} as the family table: one row per
        orthogroup, one column per species, the genes of that species listed in
        the cell. This is what \code{read\_orthofinder\_file} parses.}
  \step{Decide which genes are the \emph{orthologs}, which is a separate
        question --- see the pitfall below. OrthoFinder answers it in
        \file{Orthologues/}, whose pairwise files mark each relation as
        one-to-one, one-to-many or many-to-many; the genes appearing alone on
        both sides of a row are the unambiguous one-to-one orthologs.}
  \step{Check the result against your expression table before using it: family
        sizes, how many species each family covers, and above all how many of
        your genes actually landed in a family.}
\end{workflow}

\examplescript{python/examples/gene_family_detection.py}{%
  Surveys an \file{Orthogroups.tsv}, checks it against an expression table, and
  derives the one-to-one ortholog list from an \file{Orthologues/} directory.
  It covers the handoff only --- OrthoFinder itself is not run here.}

\begin{codefragment}[language=bash,caption={The external step, and what to keep from it}]
# one protein per gene: isoforms inflate every family
# (use your annotation's primary-transcript flag)

orthofinder -f proteomes/ -t 16

# what Tensor Omics reads afterwards:
#   Orthogroups/Orthogroups.tsv
#       the FAMILIES: row = orthogroup, column = species
#   Orthologues/<species>/<a>__v__<b>.tsv
#       the ORTHOLOGY relations: 3 columns, Orthogroup + one
#       per species; a row with ONE gene on each side is a
#       one-to-one ortholog pair
#   Orthogroups/Orthogroups_UnassignedGenes.tsv
#       genes in no family -- they get the sentinel 0 in TOX
\end{codefragment}

\begin{parameters}
  \param{primary transcripts}{Whether one or many proteins represent each
    gene}{Always reduce first. Every isoform kept becomes an apparent paralog
    and biases family size, centroids and every divergence test built on them}
  \param{species set}{Which genomes define the families}{Adding a distant
    species merges families; removing one can split them. The set is part of
    the result, not a free parameter --- record it}
  \param{ortholog rule}{Which family members count as the conserved reference}{%
    The one-to-one relations in \file{Orthologues/} are the strict choice.
    The single-copy orthogroups file is stricter still, but excludes exactly the
    duplicated families a divergence analysis is about. State which you used}
  \param{inflation (MCL, \optn{-I})}{Granularity, if you refine large
    orthogroups further}{Higher (\code{2.0}--\code{4.0}) yields smaller, tighter
    families; lower yields larger, more inclusive ones}
\end{parameters}

\begin{expectedresults}
  The bundled \file{material/Orthogroups.tsv} holds 15\,512 orthogroups over
  four mammal species, family sizes from 2 to 143 with a median of 4. Of these
  12\,842 contain all four species and 2\,670 are missing at least one, while
  4\,004 hold more genes than species and therefore contain duplications ---
  those are the families an SND analysis has anything to say about. Against the
  bundled expression table, 81\,960 of 88\,327 genes land in a family and
  6\,367 carry the unassigned sentinel. A far larger unassigned fraction almost
  always means the identifiers do not match, not that the genes are orphans.
\end{expectedresults}

\begin{pitfall}[Pitfall: taking an orthogroup for a set of orthologs]
  They are different things, and Tensor Omics needs both. An \emph{orthogroup}
  is the set of genes descended from a single gene in the last common ancestor
  of \emph{all} the species; \emph{orthologues} are genes descended from a
  single gene in the ancestor of a \emph{pair} of species; \emph{paralogues}
  arose by duplication within an orthogroup. \file{Orthogroups.tsv} therefore
  gives you families that mix orthologs and paralogs, and reading it does not
  tell you which is which. \code{group\_centroid\_orthologs} asks for exactly
  that distinction, and it comes from \file{Orthologues/}, not from the
  orthogroup table.
\end{pitfall}

\begin{pitfall}[Pitfall: retreating to the single-copy orthogroups]
  \file{Orthogroups\_SingleCopyOrthologues.txt} lists the families with exactly
  one gene per species, which makes the ortholog question trivial --- every
  member is one. It is also useless for the analyses in this manual: a family
  with no duplicated copies has no paralog to test for subfunctionalization, no
  outlier to detect against its siblings. Use it to sanity-check a pipeline,
  never as the working family set.
\end{pitfall}

\begin{pitfall}[Pitfall: identifiers that nearly match]
  The family table and the expression table are joined by exact string
  equality. A version suffix present in one and stripped in the other, or a
  transcript identifier against a gene identifier, produces no error anywhere
  --- the genes simply come back unassigned and quietly leave the analysis. The
  example script reports the assigned count for this reason; check it against
  what you expect before going on.
\end{pitfall}

\begin{pitfall}[Pitfall: isoforms as paralogs]
  Several splice isoforms of one gene look exactly like several copies of it.
  They inflate family sizes, drag centroids towards whichever gene has the most
  isoforms, and can present as a dosage effect that is purely an annotation
  artefact. Reduce to primary transcripts before OrthoFinder, not after.
\end{pitfall}

\begin{interpretation}
  The family grouping is not a preprocessing detail; it defines the comparison.
  Every later recipe asks how a gene differs from \emph{its family}, so a
  family that is too inclusive hides divergence inside itself, and one that is
  too fragmented has nothing to compare against. When a result looks surprising,
  the family it came from is the first thing to inspect.

  The ortholog designation carries a second, independent decision: what counts
  as the ancestral state. A centroid over one-to-one orthologs is a claim about
  conserved expression, and the divergence measured against it is divergence
  from that conserved state. A centroid over all members measures divergence
  from the family as it is now, duplications included. Both are legitimate;
  they answer different questions, and no later recipe can recover which one
  you meant. Record the species set, the ortholog rule and the OrthoFinder
  version alongside the results.
\end{interpretation}

\begin{references}
  \reference{Emms \& Kelly (2019). OrthoFinder: phylogenetic orthology
    inference for comparative genomics. \emph{Genome Biology} 20:238.}
  \reference{Emms \& Kelly (2015). OrthoFinder: solving fundamental biases in
    whole genome comparisons dramatically improves orthogroup inference
    accuracy. \emph{Genome Biology} 16:157.}
\end{references}

\end{recipe}
